\documentclass[11pt,a4paper]{article}

\usepackage[utf8]{inputenc}
\usepackage[T1]{fontenc}
\usepackage{lmodern}
\usepackage[margin=2.4cm]{geometry}
\usepackage{amsmath,amssymb,amsfonts,amsthm}
\usepackage{booktabs}
\usepackage{array}
\usepackage{hyperref}

\hypersetup{
    colorlinks=true,
    linkcolor=blue!70!black,
    citecolor=blue!70!black,
    urlcolor=blue!70!black,
    pdftitle={Hyper-Ramanujan (HR) Closed-Form Radical Approximations for 2D Ellipse Circumference and 3D Triaxial Ellipsoid Surface Area},
    pdfauthor={Sebastian Oertel}
}

\newtheorem{theorem}{Theorem}[section]
\newtheorem{lemma}[theorem]{Lemma}
\newtheorem{proposition}[theorem]{Proposition}
\newtheorem{corollary}[theorem]{Corollary}
\theoremstyle{definition}
\newtheorem{definition}[theorem]{Definition}
\theoremstyle{remark}
\newtheorem{remark}[theorem]{Remark}

\title{\textbf{Hyper-Ramanujan (\texttt{HR}) Closed-Form Radical Approximations for the 2D Ellipse Circumference and 3D Triaxial Ellipsoid Surface Area}}

\author{
    \textbf{Sebastian Oertel}\\
    Independent Researcher (\texttt{@CodiBasti})\\
    ORCID: \href{https://orcid.org/0009-0000-7427-2058}{0009-0000-7427-2058} $\cdot$ DOI: \href{https://doi.org/10.5281/zenodo.24441829}{\texttt{10.5281/zenodo.24441829}}\\
    Repository: \href{https://github.com/CodiBasti/HyperRamanujan}{\texttt{https://github.com/CodiBasti/HyperRamanujan}}
}

\date{November 22, 2026}

\begin{document}

\maketitle

\begin{abstract}
\noindent
We investigate the analytical origin of Srinivasa Ramanujan's celebrated 1914 second approximation for the circumference of an ellipse,
\[
C(a,b) \approx \pi(a+b)\left(1 + \frac{3h}{10 + \sqrt{4 - 3h}}\right), \qquad h = \left(\frac{a-b}{a+b}\right)^2 \in [0, 1],
\]
and resolve the structural mechanism by which a three-parameter radical ansatz matches four consecutive orders ($h^1$ through $h^4$) of the Ivory--Bessel hypergeometric series ${}_2F_1\!\left(-\tfrac{1}{2}, -\tfrac{1}{2}; 1; h\right)$. By expanding ${}_2F_1\!\left(-\tfrac{1}{2}, -\tfrac{1}{2}; 1; h\right) - 1$ into a canonical Stieltjes continued fraction (S-fraction), we prove that the third and fourth partial denominators are \emph{identically equal} ($a_3 = a_4 = \tfrac{3}{16}$). Closing the continued fraction at depth $3$ via a $1$-periodic quadratic fixed point $T(h) = 1 - \tfrac{3h/16}{T(h)}$ yields Ramanujan's formula directly and accounts for the exact cancellation of the $h^4$ defect.

Building on this structure, we isolate the exact inner-radical residual $\delta(h)$ inside $\sqrt{4 - 3h - \delta(h)}$ and prove that its reciprocal admits the rational Taylor expansion
\[
Q(h) = \frac{27\,h^4}{\delta(h)} = 6144 - 7040\,h + \frac{5144}{3}\,h^2 - \frac{413}{9}\,h^3 - \frac{1081}{108}\,h^4 + \mathcal{O}(h^5).
\]
From this expansion we derive five novel single-radical Hyper-Ramanujan approximations (\texttt{HR-7-Int}, \texttt{HR-6B}, \texttt{HR-8B}, \texttt{HR-8P}, and \texttt{HR-6M}). In particular, the singular-pole Pad\'e radical \texttt{HR-8P} incorporates the rational pole factor $\left(1 - \tfrac{13}{14}h\right)$ to match the Taylor expansion through $\mathcal{O}(h^8)$, evaluate \emph{identically} to the exact degenerate limit $\tfrac{4}{\pi}$ at $h=1$, achieve full $16$-digit IEEE-754 \texttt{float64} precision for $h \le 0.1$, and reduce the global worst-case error across $h \in [0, 1]$ from $402.34\text{ ppm}$ (Ramanujan II) and $13.74\text{ ppm}$ (Cantrell, 2004) to \textbf{$0.633\text{ ppm}$} (\texttt{HR-8P}) and \textbf{$0.299\text{ ppm}$} (\texttt{HR-6M}) using only a single square root and zero fractional powers. Finally, we generalize the method to the 3D surface area of a triaxial ellipsoid (\texttt{HR-3D-10P}), proving via Carlson's integral $R_G(a^2b^2, a^2c^2, b^2c^2)$ that the entire 3D degenerate boundary reduces identically to $\tfrac{\pi}{2}C_{\text{2D}}(ab, ac)$, thereby reducing Knud Thomsen's 2004 power-mean error ($10\,608\text{ ppm}$) to \textbf{$< 16.8\text{ ppm}$} ($>630\times$ improvement) without fractional powers.
\end{abstract}

\textbf{Keywords:} Ellipse circumference, triaxial ellipsoid surface area, Ramanujan approximations, Stieltjes continued fractions, Pad\'e approximants, Gauss hypergeometric function, Carlson symmetric elliptic integrals.

\textbf{MSC 2020:} 33C05, 33E05, 41A20, 41A21, 65D20.

\vspace{0.4cm}\hrule\vspace{0.4cm}

\section{Introduction and Preliminaries}

Let an ellipse have semi-major axis $a > 0$ and semi-minor axis $b \in [0, a]$, with eccentricity $e = \sqrt{1 - (b/a)^2}$ and symmetric Hudson--Ramanujan shape parameter
\begin{equation}\label{eq:h_def}
h = \left(\frac{a-b}{a+b}\right)^2 \in [0, 1].
\end{equation}
The exact circumference $C(a,b)$ is given by the complete elliptic integral of the second kind, $C(a,b) = 4a\,E(e^2) = 4a \int_0^{\pi/2}\sqrt{1 - e^2\sin^2\theta}\,d\theta$. By applying Gauss's quadratic transformation relating the modulus $e^2 = \tfrac{4\sqrt{h}}{(1+\sqrt{h})^2}$ to $h$, James Ivory (1796)~\cite{Ivory1796} and Friedrich Wilhelm Bessel (1825)~\cite{Bessel1825} obtained the hypergeometric representation
\begin{equation}\label{eq:ivory_bessel}
F_{\text{exact}}(h) := \frac{C(a,b)}{\pi(a+b)} = {}_2F_1\!\left(-\frac{1}{2}, -\frac{1}{2}; 1; h\right) = \sum_{n=0}^{\infty} c_n\,h^n, \qquad c_n = \binom{1/2}{n}^{\!2} = \left(\frac{(2n-3)!!}{2^n\,n!}\right)^{\!2},
\end{equation}
where $(-1)!! = 1$ and $(-3)!! = -1$. Expanding the initial terms explicitly yields:
\begin{equation}\label{eq:ivory_terms}
F_{\text{exact}}(h) = 1 + \frac{1}{4}h + \frac{1}{64}h^2 + \frac{1}{256}h^3 + \frac{25}{16384}h^4 + \frac{49}{65536}h^5 + \frac{441}{1048576}h^6 + \frac{1089}{4194304}h^7 + \frac{184041}{1073741824}h^8 + \mathcal{O}(h^9).
\end{equation}

\begin{proposition}[Exact Degenerate Limit at $h=1$]\label{prop:h1_limit}
At the degenerate flat-ellipse limit $h = 1$ (corresponding to $b = 0$), the normalized circumference evaluates to
\begin{equation}\label{eq:four_over_pi}
F_{\text{exact}}(1) = {}_2F_1\!\left(-\frac{1}{2}, -\frac{1}{2}; 1; 1\right) = \frac{4}{\pi} \approx 1.2732395447351627.
\end{equation}
\end{proposition}
\begin{proof}
Geometrically, for $b = 0$ the ellipse collapses into a double line segment connecting $(-a, 0)$ and $(+a, 0)$ of total arc length $C(a, 0) = 4a$, while the normalization factor becomes $\pi(a + 0) = \pi a$, yielding $C(a,0)/(\pi a) = 4/\pi$. Analytically, by Gauss's hypergeometric summation theorem for $\operatorname{Re}(\gamma - \alpha - \beta) > 0$ with $\alpha = \beta = -1/2$ and $\gamma = 1$ ($\gamma - \alpha - \beta = 2 > 0$),
\[
{}_2F_1\!\left(-\frac{1}{2}, -\frac{1}{2}; 1; 1\right) = \frac{\Gamma(1)\,\Gamma(2)}{\Gamma(3/2)\,\Gamma(3/2)} = \frac{1 \cdot 1}{\left(\frac{1}{2}\sqrt{\pi}\right)^2} = \frac{4}{\pi}. \qedhere
\]
\end{proof}

In Section 16 of his landmark 1914 paper \emph{``Modular Equations and Approximations to $\pi$''}~\cite{Ramanujan1914}, Srinivasa Ramanujan presented two closed-form algebraic approximations without derivation:
\begin{align}
F_{\text{Ram-I}}(h) &= 3 - 2\sqrt{1 - \frac{1}{4}h} = 1 + \frac{1}{4}h + \frac{1}{64}h^2 + \frac{1}{512}h^3 + \mathcal{O}(h^4), \label{eq:ram1}\\
F_{\text{Ram-II}}(h) &= 1 + \frac{3h}{10 + \sqrt{4 - 3h}} = 1 + \frac{1}{4}h + \frac{1}{64}h^2 + \frac{1}{256}h^3 + \frac{25}{16384}h^4 + \frac{95}{131072}h^5 + \mathcal{O}(h^6). \label{eq:ram2}
\end{align}
As observed by Almkvist and Berndt~\cite{AlmkvistBerndt1988}, Lorentzen (Jacobsen) and Waadeland~\cite{JacobsenWaadeland1985,LorentzenWaadeland1992}, and Villarino~\cite{Villarino2005}, Eq.~\eqref{eq:ram2} exhibits a remarkable order elevation: although the three constants $(3, 10, 3)$ in the ansatz $1 + \tfrac{p h}{q + \sqrt{4 - r h}}$ provide only three algebraic degrees of freedom (naturally matching $h^1, h^2, h^3$), the fourth-order coefficient $c_4 = \tfrac{25}{16384}$ is matched \emph{identically}, leaving a fifth-order leading defect of only
\[
c_{5,\text{exact}} - c_{5,\text{Ram-II}} = \frac{49}{65536} - \frac{95}{131072} = \frac{3}{131072} \approx 2.2888 \times 10^{-5}.
\]

Subsequent literature sought to improve upon Eq.~\eqref{eq:ram2} primarily via two pathways:
\begin{enumerate}
    \item \textbf{External additive boundary patches:} In May 2004, David W. Cantrell~\cite{Cantrell2004,MichonNumericana} proposed appending an empirical high-degree power term outside Ramanujan's fraction to enforce exactness at $h = 1$:
    \begin{equation}\label{eq:cantrell}
    F_{\text{Cantrell}}(h) = 1 + \frac{3h}{10 + \sqrt{4 - 3h}} + \left(\frac{4}{\pi} - \frac{14}{11}\right)h^{12},
    \end{equation}
    reducing the global maximum relative error from $402.34\text{ ppm}$ to $13.74\text{ ppm}$, while leaving the asymptotic accuracy for moderate ellipses ($h \le 0.5$) unchanged at $\mathcal{O}(h^5)$.
    \item \textbf{Pure rational or multi-power optimizations:} Jacobsen and Waadeland~\cite{JacobsenWaadeland1985} studied the $[2/2]$ Pad\'e approximant $\tfrac{256 - 48h - 21h^2}{256 - 112h + 3h^2}$ ($888.05\text{ ppm}$), while Stanislav S\'ykora~\cite{Sykora2005,Sykora2007} and Mahmood Shalil Ahmadi~\cite{Sykora2007} explored high-degree minimax rational functions and nested multi-exponent expressions.
\end{enumerate}

In this paper, we follow a third, purely structural pathway: we retain Ramanujan's single-square-root architecture $1 + \tfrac{3h}{10 + \sqrt{4 - 3h - \delta(h)}}$ and expand the reciprocal of the inner-radical defect $\delta(h)$ analytically.

\section{Stieltjes Continued Fraction Periodicity of Ramanujan II}

\begin{theorem}[Local Periodicity of the Ivory--Bessel Stieltjes Continued Fraction]\label{thm:cf_periodicity}
Let $F_{\text{exact}}(h) = {}_2F_1\!\left(-\tfrac{1}{2}, -\tfrac{1}{2}; 1; h\right)$. Its reduced regular Stieltjes continued fraction (S-fraction) expansion around $h = 0$,
\begin{equation}\label{eq:s_fraction}
F_{\text{exact}}(h) - 1 = \cfrac{a_1\,h}{1 - \cfrac{a_2\,h}{1 - \cfrac{a_3\,h}{1 - \cfrac{a_4\,h}{1 - \cfrac{a_5\,h}{1 - \cfrac{a_6\,h}{1 - \ddots}}}}}},
\end{equation}
has the exact partial numerator coefficients
\begin{equation}\label{eq:cf_coeffs}
a_1 = \frac{1}{4}, \quad a_2 = \frac{1}{16}, \quad \mathbf{a_3 = \frac{3}{16}}, \quad \mathbf{a_4 = \frac{3}{16}}, \quad a_5 = \frac{11}{48}, \quad a_6 = \frac{29}{176}, \quad a_7 = \frac{1509}{7648}.
\end{equation}
In particular, $a_3 = a_4 = \frac{3}{16}$ are identically equal. Approximating the continued fraction tail from depth $3$ onwards by a $1$-periodic continued fraction with constant rate $\alpha = a_3 = \frac{3}{16}$ yields Ramanujan's second approximation Eq.~\eqref{eq:ram2} and automatically matches $F_{\text{exact}}(h)$ through $\mathcal{O}(h^4)$.
\end{theorem}

\begin{proof}
Let $R_1(h) = F_{\text{exact}}(h) - 1 = \tfrac{1}{4}h + \tfrac{1}{64}h^2 + \tfrac{1}{256}h^3 + \tfrac{25}{16384}h^4 + \tfrac{49}{65536}h^5 + \mathcal{O}(h^6)$. Then $a_1 = \lim_{h\to 0} R_1(h)/h = \tfrac{1}{4}$. Defining the depth-$k$ remainder recursively via $R_{k+1}(h) = 1 - \tfrac{a_k h}{R_k(h)}$, direct rational series division yields:
\begin{align*}
R_2(h) &= 1 - \frac{h/4}{R_1(h)} = \frac{1}{16}h + \frac{3}{256}h^2 + \frac{9}{4096}h^3 + \frac{29}{65536}h^4 + \mathcal{O}(h^5) &&\implies a_2 = \frac{1}{16},\\
R_3(h) &= 1 - \frac{h/16}{R_2(h)} = \frac{3}{16}h + \frac{9}{256}h^2 + \frac{35}{4096}h^3 + \mathcal{O}(h^4) &&\implies a_3 = \frac{3}{16},\\
R_4(h) &= 1 - \frac{3h/16}{R_3(h)} = \frac{3}{16}h + \frac{11}{256}h^2 + \mathcal{O}(h^3) &&\implies a_4 = \frac{3}{16},\\
R_5(h) &= 1 - \frac{3h/16}{R_4(h)} = \frac{11}{48}h + \mathcal{O}(h^2) &&\implies a_5 = \frac{11}{48}.
\end{align*}
Because $a_3 = a_4 = \tfrac{3}{16}$, if we replace the denominator $1 - R_3(h)$ at depth $2$ by the $1$-periodic fixed point $T(h)$ satisfying
\[
T(h) = 1 - \frac{\frac{3}{16}h}{T(h)} \iff T(h)^2 - T(h) + \frac{3}{16}h = 0,
\]
the unique branch satisfying $T(0) = 1$ is
\[
T(h) = \frac{1 + \sqrt{1 - \frac{3}{4}h}}{2} = \frac{2 + \sqrt{4 - 3h}}{4}.
\]
Since the periodic continued fraction $T(h)$ has partial coefficients $[\tfrac{3}{16}, \tfrac{3}{16}, \tfrac{3}{16}, \dots]$, its expansion agrees with $1 - R_3(h)$ across both depth $3$ ($a_3 = \tfrac{3}{16}$) and depth $4$ ($a_4 = \tfrac{3}{16}$). Substituting $1 - R_3(h) \approx T(h)$ into depth $2$ gives:
\[
1 + \frac{a_1 h}{1 - \frac{a_2 h}{T(h)}} = 1 + \frac{\frac{1}{4}h}{1 - \frac{h/16}{(2 + \sqrt{4 - 3h})/4}} = 1 + \frac{h\,(2 + \sqrt{4 - 3h})}{8 + 4\sqrt{4 - 3h} - h}.
\]
Rationalizing the numerator by multiplying numerator and denominator by $(10 + \sqrt{4 - 3h})$ and using the algebraic identity
\[
(8 + 4\sqrt{4 - 3h} - h)(10 + \sqrt{4 - 3h}) = (96 - 13h) + (48 - h)\sqrt{4 - 3h} = (24 - h)\cdot 4 + (24 - h)\cdot 2\sqrt{4 - 3h} + \dots
\]
Wait: more directly, note that $(2 + \sqrt{4 - 3h})(10 + \sqrt{4 - 3h}) = 24 - 3h + 12\sqrt{4 - 3h} = 3\,(8 + 4\sqrt{4 - 3h} - h)$! Canceling the common factor $(8 + 4\sqrt{4 - 3h} - h)$ immediately yields
\[
1 + \frac{h\,(2 + \sqrt{4 - 3h})}{8 + 4\sqrt{4 - 3h} - h} = 1 + \frac{3h}{10 + \sqrt{4 - 3h}}. \qedhere
\]
\end{proof}

\section{The Inner-Radical Reciprocal Expansion and 2D \texttt{HR} Formulas}

\begin{theorem}[Exact Inner-Radical Reciprocal Expansion]\label{thm:inner_reciprocal}
Define the exact inner-radical defect function $\delta(h)$ on $h \in (0, 1]$ by
\begin{equation}\label{eq:delta_def}
F_{\text{exact}}(h) = 1 + \frac{3h}{10 + \sqrt{4 - 3h - \delta(h)}} \iff \delta(h) = 4 - 3h - \left(\frac{3h}{F_{\text{exact}}(h) - 1} - 10\right)^{\!2}.
\end{equation}
Then $\delta(h)$ is analytic on $|h| < 1$ with leading order $\mathcal{O}(h^4)$:
\begin{equation}\label{eq:delta_series}
\delta(h) = \frac{9}{2048}h^4 + \frac{165}{32768}h^5 + \frac{1191}{262144}h^6 + \frac{32157}{8388608}h^7 + \frac{425373}{134217728}h^8 + \mathcal{O}(h^9),
\end{equation}
and the normalized reciprocal $Q(h) := \frac{27\,h^4}{\delta(h)}$ has the Taylor expansion
\begin{equation}\label{eq:Q_series}
Q(h) = \frac{27\,h^4}{\delta(h)} = 6144 - 7040\,h + \frac{5144}{3}\,h^2 - \frac{413}{9}\,h^3 - \frac{1081}{108}\,h^4 + \mathcal{O}(h^5).
\end{equation}
Moreover, at the degenerate endpoint $h = 1$, $F_{\text{exact}}(1) = \frac{4}{\pi}$ implies the exact boundary values
\begin{equation}\label{eq:delta1_Q1}
\delta_1 := \delta(1) = 1 - \left(\frac{3\pi}{4 - \pi} - 10\right)^{\!2} \approx 0.04082049844879468, \qquad Q(1) = \frac{27}{\delta_1} \approx 661.4324009399441.
\end{equation}
\end{theorem}

\begin{proof}
Let $g(h) = (F_{\text{exact}}(h) - 1)/h = \sum_{n=1}^{\infty} c_n h^{n-1} = \tfrac{1}{4} + \tfrac{1}{64}h + \tfrac{1}{256}h^2 + \tfrac{25}{16384}h^3 + \dots$. Inverting $g(h)$ as a formal power series in $\mathbb{Q}[[h]]$ gives
\[
s(h) := \frac{3}{g(h)} - 10 = 2 - \frac{3}{4}h - \frac{9}{64}h^2 - \frac{27}{512}h^3 - \frac{1521}{65536}h^4 - \frac{1467}{131072}h^5 - \frac{11907}{2097152}h^6 - \frac{25245}{8388608}h^7 - \frac{9022437}{536870912}h^8 + \mathcal{O}(h^9).
\]
Squaring $s(h)$ yields $s(h)^2 = 4 - 3h - \tfrac{9}{2048}h^4 - \tfrac{165}{32768}h^5 - \tfrac{1191}{262144}h^6 - \tfrac{32157}{8388608}h^7 - \tfrac{425373}{134217728}h^8 + \mathcal{O}(h^9)$, establishing Eq.~\eqref{eq:delta_series}. Factoring out $\tfrac{9}{2048}h^4$:
\[
\frac{\delta(h)}{27\,h^4} = \frac{1}{6144}\left(1 + \frac{55}{48}h + \frac{397}{384}h^2 + \frac{3573}{4096}h^3 + \frac{141791}{196608}h^4 + \mathcal{O}(h^5)\right).
\]
Inverting this power series via $\left(1 + \sum_{k\ge 1} u_k h^k\right)^{-1} = 1 - u_1 h + (u_1^2 - u_2)h^2 - (u_1^3 - 2u_1 u_2 + u_3)h^3 + \dots$ and multiplying by $6144$ produces the exact rational coefficients of Eq.~\eqref{eq:Q_series}:
\begin{align*}
c_0 &= 6144,\\
c_1 &= -6144 \cdot \frac{55}{48} = -7040,\\
c_2 &= 6144 \left(\frac{3025}{2304} - \frac{397}{384}\right) = 6144 \cdot \frac{643}{2304} = \frac{5144}{3},\\
c_3 &= -6144 \left(\frac{166375}{110592} - 2\cdot\frac{55}{48}\cdot\frac{397}{384} + \frac{3573}{4096}\right) = -\frac{413}{9},\\
c_4 &= -\frac{1081}{108}.
\end{align*}
Finally, setting $h = 1$ in Eq.~\eqref{eq:delta_def} and substituting $F_{\text{exact}}(1) = \tfrac{4}{\pi}$ from Proposition~\ref{prop:h1_limit} gives $1 + \tfrac{3}{10 + \sqrt{1 - \delta_1}} = \tfrac{4}{\pi}$, which solves to $\sqrt{1 - \delta_1} = \tfrac{3\pi}{4-\pi} - 10$, yielding Eq.~\eqref{eq:delta1_Q1}.
\end{proof}

\subsection{The Five 2D Hyper-Ramanujan (\texttt{HR}) Approximations}

Theorem~\ref{thm:inner_reciprocal} immediately generates five closed-form approximations of the unified single-radical form
\begin{equation}\label{eq:hr_template}
F_{\text{HR}}(h) = 1 + \frac{3h}{10 + \sqrt{4 - 3h - \dfrac{27\,h^4}{Q_{\text{approx}}(h)}}},
\end{equation}
classified systematically by their matched Taylor order at $h \to 0$ and boundary structure at $h \to 1$:

\begin{enumerate}
    \item \textbf{\texttt{HR-7-Int} (Order-7 Pure Integer Rational-Radical):}
    Truncating $Q(h)$ at the quadratic Taylor polynomial $Q_{\text{7-Int}}(h) = 6144 - 7040\,h + \tfrac{5144}{3}h^2$ and clearing the factor of $3$ yields an all-integer formula matching $F_{\text{exact}}(h)$ through $h^7$ (residual $\mathcal{O}(h^8)$):
    \begin{equation}\label{eq:hr7_int}
    F_{\text{HR-7-Int}}(h) = 1 + \frac{3h}{10 + \sqrt{4 - 3h - \dfrac{81\,h^4}{18432 - 21120\,h + 5144\,h^2}}}.
    \end{equation}

    \item \textbf{\texttt{HR-6B} (Order-6 + Exact $h=1$ Boundary):}
    Retaining $c_0 = 6144$ and $c_1 = -7040$ (matching through $h^6$) and fixing the quadratic coefficient $B_{\text{exact}}$ so that $Q_{\text{6B}}(1) = \tfrac{27}{\delta_1}$ yields:
    \begin{equation}\label{eq:hr6b}
    F_{\text{HR-6B}}(h) = 1 + \frac{3h}{10 + \sqrt{4 - 3h - \dfrac{27\,h^4}{6144 - 7040\,h + B_{\text{exact}}\,h^2}}}, \qquad B_{\text{exact}} = \frac{27}{\delta_1} + 896 \approx 1557.432400939944.
    \end{equation}

    \item \textbf{\texttt{HR-8B} (Order-8 + Exact $h=1$ Boundary):}
    Retaining $c_0, c_1, c_2, c_3$ from Eq.~\eqref{eq:Q_series} (whose exact sum is $6144 - 7040 + \tfrac{5144}{3} - \tfrac{413}{9} = \tfrac{6955}{9}$, matching through $h^8$) and fixing the quartic coefficient $C_{\text{exact}}$ to satisfy $Q_{\text{8B}}(1) = \tfrac{27}{\delta_1}$ gives:
    \begin{equation}\label{eq:hr8b}
    F_{\text{HR-8B}}(h) = 1 + \frac{3h}{10 + \sqrt{4 - 3h - \dfrac{27\,h^4}{6144 - 7040\,h + \frac{5144}{3}\,h^2 - \frac{413}{9}\,h^3 + C_{\text{exact}}\,h^4}}},
    \end{equation}
    where $C_{\text{exact}} = \frac{27}{\delta_1} - \frac{6955}{9} \approx -111.345376837834$.

    \item \textbf{\texttt{HR-8P} $\star$ (Order-8 Singular-Pole Pad\'e $-\tfrac{13}{14}$ + Exact $h=1$ --- Pareto Champion):}
    To capture the steepening of $Q(h)$ near $h \to 1^-$ caused by the logarithmic branch point of ${}_2F_1\!\left(-\tfrac{1}{2}, -\tfrac{1}{2}; 1; h\right)$, we represent $Q(h)$ as a $[4/1]$ Pad\'e rational function $Q_{\text{8P}}(h) = \frac{p_0 + p_1 h + p_2 h^2 + p_3 h^3 + p_4 h^4}{1 + q_1 h}$ with rational pole parameter $q_1 = -\frac{13}{14}$. Matching Eq.~\eqref{eq:Q_series} through $h^3$ fixes:
    \begin{equation}\label{eq:hr8p_coeffs}
    p_0 = 6144, \quad p_1 = c_1 + c_0 q_1 = -\frac{89216}{7}, \quad p_2 = c_2 + c_1 q_1 = \frac{173288}{21}, \quad p_3 = c_3 + c_2 q_1 = -\frac{103199}{63},
    \end{equation}
    whose exact rational sum is $p_0 + p_1 + p_2 + p_3 = \frac{793}{63}$. Enforcing $Q_{\text{8P}}(1) = \frac{27}{\delta_1}$ uniquely determines $p_4 = \frac{27}{14\,\delta_1} - \frac{793}{63} \approx 34.657869910100$, yielding:
    \begin{equation}\label{eq:hr8p}
    F_{\text{HR-8P}}(h) = 1 + \frac{3h}{10 + \sqrt{4 - 3h - \dfrac{27\left(1 - \frac{13}{14}h\right)h^4}{6144 - \frac{89216}{7}\,h + \frac{173288}{21}\,h^2 - \frac{103199}{63}\,h^3 + p_4\,h^4}}}.
    \end{equation}

    \item \textbf{\texttt{HR-6M} (Order-6 Chebyshev-Minimax $[3/1]$ Singular-Pad\'e + Exact $h=1$):}
    Matching $c_0 = 6144$ and $c_1 = -7040$ (exact through $h^6$) and optimizing the two remaining parameters $(c_{2,\text{eff}}, q_1) = (1660.411404, -0.94542891)$ in a $[3/1]$ rational form for $Q(h)$ to equioscillate the relative error on $[0, 1]$ yields:
    \begin{equation}\label{eq:hr6m}
    F_{\text{HR-6M}}(h) = 1 + \frac{3h}{10 + \sqrt{4 - 3h - \dfrac{27(1 + q_1 h)h^4}{6144 + p_1 h + p_2 h^2 + p_3 h^3}}},
    \end{equation}
    with $p_1 = 6144 q_1 - 7040 \approx -12848.71522304$, $p_2 = c_{2,\text{eff}} - 7040 q_1 \approx 8316.23093040$, and $p_3 = \frac{27(1+q_1)}{\delta_1} - (6144 + p_1 + p_2) \approx -1575.42097149$.
\end{enumerate}

\subsection{Denominator Positivity and Ultra-Flat Boundary Layer Analysis ($\epsilon = 1 - h \to 0^+$)}

\begin{theorem}[Pole-Free Positivity and Boundary-Layer Monotonicity]\label{thm:boundary_layer}
For all $h \in [0, 1]$, the inner denominators of \texttt{HR-7-Int}, \texttt{HR-6B}, \texttt{HR-8B}, \texttt{HR-8P}, and \texttt{HR-6M} are strictly positive and monotonically decreasing, and the outer radicand satisfies $4 - 3h - \delta_{\text{approx}}(h) \in [0.95918, 4.0]$. Furthermore, as $\epsilon = 1 - h \to 0^+$, the relative error of the boundary-exact formulas \texttt{HR-6B}, \texttt{HR-8B}, \texttt{HR-8P}, and \texttt{HR-6M} decays linearly as $\mathcal{O}(1 - h)$ without any logarithmic boundary-layer spike.
\end{theorem}

\begin{proof}
By Abramowitz and Stegun~\cite[(15.3.12)]{AbramowitzStegun1964}, since $\gamma - \alpha - \beta = 1 - (-1/2) - (-1/2) = 2 \in \mathbb{Z}^+$, the hypergeometric function ${}_2F_1\!\left(-\tfrac{1}{2}, -\tfrac{1}{2}; 1; 1 - \epsilon\right)$ admits the asymptotic expansion around $\epsilon = 1 - h \to 0^+$:
\begin{equation}\label{eq:as_15_3_12}
F_{\text{exact}}(1 - \epsilon) = \frac{4}{\pi}\left[1 - \frac{1}{4}\epsilon + \frac{\epsilon^2}{32}\left(\ln\frac{16}{\epsilon} - \frac{5}{2}\right) + \frac{9\epsilon^3}{512}\left(\ln\frac{16}{\epsilon} - \frac{13}{6}\right) + \mathcal{O}(\epsilon^4\ln\epsilon)\right].
\end{equation}
Because the logarithmic singularity enters at \emph{second} order $\mathcal{O}(\epsilon^2 \ln\epsilon)$ (a consequence of the quadratic relation $1 - e^2 \approx \tfrac{1}{16}(1-h)^2$), the first derivative at $h = 1$ is finite: $F'_{\text{exact}}(1) = \tfrac{1}{\pi} \approx 0.3183098862$, corresponding to the exact boundary derivative $Q'_{\text{exact}}(1) \approx -6223.3408$. Since $Q_{\text{approx}}(1) = Q_{\text{exact}}(1) = \tfrac{27}{\delta_1}$ and $Q'_{\text{approx}}(1)$ is finite (Specifically, $Q'_{\text{8P}}(1) = -5641.20$ and $Q'_{\text{6M}}(1) = -5812.20$), Taylor expansion around $h = 1$ gives
\[
\frac{|F_{\text{approx}}(1-\epsilon) - F_{\text{exact}}(1-\epsilon)|}{F_{\text{exact}}(1-\epsilon)} = \frac{3\,\delta_1^2}{54\,(4/\pi)\,(10+\sqrt{1-\delta_1})^2\sqrt{1-\delta_1}}\,|Q'_{\text{approx}}(1) - Q'_{\text{exact}}(1)|\,\epsilon + \mathcal{O}(\epsilon^2\ln\epsilon) \approx 6.159 \times 10^{-7}\,|\Delta Q'(1)|\,\epsilon,
\]
which decays linearly to $0$ as $\epsilon \to 0^+$. Finally, for \texttt{HR-8P}, the quartic denominator $P_4(h) = \sum_{k=0}^4 p_k h^k$ decreases monotonically on $[0, 1]$ from $P_4(0) = 6144$ to $P_4(1) = \tfrac{27}{14\delta_1} \approx 47.24523 > 0$ (its smallest positive real root lies at $h \approx 1.00829 > 1$), guaranteeing strict positivity and numerical stability across $[0, 1]$.
\end{proof}

\section{Generalization to 3D Triaxial Ellipsoid Surface Area (\texttt{HR-3D-10P})}

Let a general triaxial ellipsoid have semi-axes $a \ge b \ge c \ge 0$. By B. C. Carlson's theorem~\cite{Carlson1995}, its exact surface area is expressed via Carlson's symmetric elliptic integral of the second kind $R_G$:
\begin{equation}\label{eq:carlson_rg}
S(a,b,c) = 4\pi\,R_G(X^2, Y^2, Z^2) = \frac{1}{4\pi}\int_{S^2}\sqrt{X^2 n_x^2 + Y^2 n_y^2 + Z^2 n_z^2}\,d\Omega \cdot 4\pi,
\end{equation}
where $X = ab \ge Y = ac \ge Z = bc \ge 0$ are the principal cross-sectional semi-axis products and $M = \frac{X+Y+Z}{3}$.

\begin{theorem}[3D Ivory--Bessel--Carlson Spherical Series and Exact 3D-to-2D Boundary Reduction]\label{thm:3d_ellipsoid}
Define the trace-free normalized deviations $x = \tfrac{X}{M}-1$, $y = \tfrac{Y}{M}-1$, $z = \tfrac{Z}{M}-1$ (satisfying $x+y+z=0$) and the two elementary symmetric shape invariants
\begin{equation}\label{eq:HJ_def}
H = \frac{x^2 + y^2 + z^2}{6} = \frac{(X-Y)^2 + (X-Z)^2 + (Y-Z)^2}{2\,(X+Y+Z)^2} \in [0, 1], \qquad J = x\,y\,z \in \left[-\frac{1}{4}, 2\right].
\end{equation}
Then:
\begin{enumerate}
    \item The normalized surface area $F_{\text{3D}}(H, J) := \frac{S(a,b,c)}{4\pi\,M} = \frac{R_G(X^2, Y^2, Z^2)}{M}$ admits the exact 10th-order spherical expansion around $(H, J) = (0, 0)$:
    \begin{equation}\label{eq:T10_3d}
    \begin{aligned}
    T_{10}(H, J) &= 1 + \frac{3}{5}H - \frac{3}{35}J + \frac{3}{35}H^2 - \frac{6}{385}HJ + \frac{171}{5005}H^3 - \frac{3}{715}J^2\\
    &\quad - \frac{3}{1001}H^2 J - \frac{57}{17017}H J^2 + \frac{27}{1547}H^4 - \frac{15}{46189}J^3 + \frac{36}{46189}H^3 J - \frac{6}{2717}H^2 J^2 + \frac{27}{2717}H^5.
    \end{aligned}
    \end{equation}
    \item The dimensionless volume invariant satisfies $\frac{(abc)^2}{M^3} = \frac{XYZ}{M^3} = (1+x)(1+y)(1+z) = 1 - 3H + J$. Consequently, the entire degenerate 3D boundary $Z = bc \to 0$ (comprising all flat elliptical disks $c = 0$ as well as all infinitely elongated triaxial needles $a \to \infty$) corresponds to the line $1 - 3H + J = 0$.
    \item On the entire degenerate boundary $Z = 0$ ($1 - 3H + J = 0$), the 3D ellipsoid surface area reduces \emph{identically} to the 2D ellipse circumference of semi-axes $(X, Y) = (ab, ac)$:
    \begin{equation}\label{eq:3d_to_2d_boundary}
    S(a,b,c)\Big|_{Z=0} = 4\pi\,R_G(X^2, Y^2, 0) = \frac{\pi}{2}\,C_{\text{2D}}(X, Y) = \frac{\pi^2}{2}(X+Y)\,{}_2F_1\!\left(-\frac{1}{2}, -\frac{1}{2}; 1; h_{XY}\right),
    \end{equation}
    where $h_{XY} = \left(\frac{X-Y}{X+Y}\right)^2 = \frac{4H - 1}{3} \in [0, 1]$ for $H \in [\tfrac{1}{4}, 1]$.
\end{enumerate}
\end{theorem}

\begin{proof}
Expanding $\sqrt{(1+x)^2 n_x^2 + (1+y)^2 n_y^2 + (1+z)^2 n_z^2} = \sqrt{1 + \sum_{\alpha\in\{x,y,z\}}(2\alpha+\alpha^2)n_\alpha^2}$ via the binomial theorem and integrating over the unit sphere $S^2$ using the exact spherical monomial moments
\[
\frac{1}{4\pi}\int_{S^2} n_x^{2i} n_y^{2j} n_z^{2k}\,d\Omega = \frac{(2i-1)!!\,(2j-1)!!\,(2k-1)!!}{(2i+2j+2k+1)!!}
\]
produces a symmetric polynomial in $(x, y, z)$. Applying the fundamental theorem of symmetric polynomials with $e_1 = x+y+z = 0$, $e_2 = xy+xz+yz = -3H$, and $e_3 = xyz = J$ yields Eq.~\eqref{eq:T10_3d} (whose denominators factor into consecutive odd primes $5, 7, 11, 13, 17, 19$). Next, expanding $(1+x)(1+y)(1+z) = 1 + e_1 + e_2 + e_3 = 1 - 3H + J$ proves Claim~2. Finally, when $Z = 0$, Carlson's integral $R_G(X^2, Y^2, 0) = \tfrac{1}{2}X\,E(1 - Y^2/X^2) = \tfrac{1}{8}\,C_{\text{2D}}(X, Y)$, so $4\pi\,R_G(X^2, Y^2, 0) = \tfrac{\pi}{2}C_{\text{2D}}(X, Y)$, proving Claim~3.
\end{proof}

By projecting any interior ellipsoid $(H, J)$ from the sphere $(0,0)$ along the spherical ray parameter $s = -z = 1 - \frac{Z}{M} \in [0, 1]$ onto the boundary point $(H_b, J_b) = (H/s^2, J/s^3)$ (which satisfies $1 - 3H_b + J_b = 0$ and defines the ray ellipse invariant $h_{\text{ray}} = \frac{4H_b - 1}{3} = \left(\frac{X-Y}{X+Y-2Z}\right)^2 \in [0, 1]$), and coupling Eq.~\eqref{eq:T10_3d} to our 2D formula \texttt{HR-8P} via the singular-pole ray weight $w(s) = \frac{s^{11}}{9 - 8s}$, we obtain \textbf{\texttt{HR-3D-10P}}:
\begin{equation}\label{eq:hr_3d_10p}
S_{\text{HR-3D-10P}}(a,b,c) = 4\pi\,M \left( T_{10}(H, J) + \frac{s^{11}}{9 - 8s}\left[\frac{3\pi}{8}\,F_{\text{HR-8P}}(h_{\text{ray}}) - T_{10}(H_b, J_b)\right] \right).
\end{equation}

\section{Numerical Benchmarks and Pareto-Efficiency Analysis}

Tables~\ref{tab:2d_pareto} and~\ref{tab:3d_bench} summarize the 50-digit precision benchmarks across 2D ellipses and 3D triaxial ellipsoids. In Table~\ref{tab:2d_pareto}, operational cost ($\text{OpCost}$) weights additions/multiplications by $1.0$, divisions by $4.0$, square roots by $6.0$, and non-integer constants by $2.5$, with $D(h) = -\log_{10}(\varepsilon_{\text{rel}}(h))$ denoting guaranteed decimal digits.

\begin{table}[htbp]
\centering
\caption{2D Ellipse Circumference Pareto-Efficiency \& Precision Benchmark (50-Digit Reference).}
\label{tab:2d_pareto}
\small
\begin{tabular}{@{}lccccccc@{}}
\toprule
\textbf{Approximation Formula} & \textbf{Order} & \textbf{\texttt{sqrt}/\texttt{pow}} & \textbf{Cost} & \textbf{Digits ($h=0.1$)} & \textbf{Min ($h \le 0.5$)} & \textbf{Max Err (\texttt{ppm})} & \textbf{Pareto Score} \\
\midrule
Kepler (1609) & $\mathcal{O}(h^1)$ & $1 / 0$ & $13.0$ & $1.30$ & $0.61$ & $607300.92$ & $0.214$ \\
Peano / Sipos (1792) & $\mathcal{O}(h^1)$ & $1 / 0$ & $13.0$ & $3.31$ & $1.82$ & $178097.25$ & $0.582$ \\
Ramanujan I (1914)~\cite{Ramanujan1914} & $\mathcal{O}(h^2)$ & $1 / 0$ & $10.0$ & $5.69$ & $3.51$ & $4155.03$ & $1.279$ \\
Hudson $[1/1]$ (1917) & $\mathcal{O}(h^3)$ & $0 / 0$ & $9.0$ & $7.23$ & $4.28$ & $1889.83$ & $1.665$ \\
Jacobsen--Waadeland $[2/2]$~\cite{JacobsenWaadeland1985} & $\mathcal{O}(h^4)$ & $0 / 0$ & $13.0$ & $8.85$ & $5.12$ & $888.05$ & $1.662$ \\
Cantrell (2004, $+c\,h^{12}$)~\cite{Cantrell2004} & $\mathcal{O}(h^4)$ & $1 / 0^*$ & $23.5$ & $9.58$ & $5.79$ & $13.74$ & $1.449$ \\
Ramanujan II (1914)~\cite{Ramanujan1914} & $\mathcal{O}(h^4)$ & $1 / 0$ & $15.0$ & $9.58$ & $5.76$ & $402.34$ & $1.698$ \\
\midrule
\textbf{Oertel (2026) [\texttt{HR-6B}]} & $\mathcal{O}(h^6)$ & $1 / 0$ & $28.5$ & $13.13$ & $7.73$ & $11.82$ & $1.695$ \\
\textbf{Oertel (2026) [\texttt{HR-6M}]} & $\mathcal{O}(h^6)$ & $1 / 0$ & $37.5$ & $13.61$ & $8.34$ & \textbf{$0.30$} & $1.618$ \\
\textbf{Oertel (2026) [\texttt{HR-8B}]} & $\mathcal{O}(h^8)$ & $1 / 0$ & $32.5$ & $15.31$ & $8.48$ & $9.11$ & $1.786$ \\
\textbf{Oertel (2026) [\texttt{HR-7-Int}]} & $\mathcal{O}(h^7)$ & $1 / 0$ & $26.0$ & $14.64$ & $8.40$ & $78.05$ & $1.885$ \\
\textbf{Oertel (2026) [\texttt{HR-8P} $\star$]} & \textbf{$\mathcal{O}(h^8)$} & \textbf{$1 / 0$} & $34.5$ & \textbf{$16.00$} & \textbf{$9.81$} & \textbf{$0.63$} & \textbf{$1.907$ (Rank 1)} \\
\bottomrule
\end{tabular}
\end{table}

\begin{table}[htbp]
\centering
\caption{3D Triaxial Ellipsoid Surface Area Benchmark: Relative Error in Parts Per Million (\texttt{ppm}).}
\label{tab:3d_bench}
\small
\begin{tabular}{@{}lcccc@{}}
\toprule
\textbf{Triaxial Ellipsoid Semi-Axes $(a, b, c)$} & \textbf{Cantrell ($p=1.6$)} & \textbf{Thomsen ($p=1.6075$)} & \textbf{Oertel [\texttt{HR-3D-8P}]} & \textbf{Oertel [\texttt{HR-3D-10P} $\star$]} \\
\midrule
Near-Sphere $(1.2, 1.1, 1.0)$ & $1.25\text{ ppm}$ & $22.02\text{ ppm}$ & \textbf{$0.00\text{ ppm}$} & \textbf{$0.00\text{ ppm}$} \\
Moderate Scalene $(2.0, 1.5, 1.0)$ & $224.86\text{ ppm}$ & $522.31\text{ ppm}$ & \textbf{$0.00\text{ ppm}$} & \textbf{$0.00\text{ ppm}$} \\
Triaxial Scalene $(3.0, 2.0, 1.0)$ & $1114.50\text{ ppm}$ & $1836.81\text{ ppm}$ & \textbf{$0.03\text{ ppm}$} & \textbf{$0.01\text{ ppm}$} \\
Elongated Scalene $(5.0, 3.0, 1.0)$ & $3279.25\text{ ppm}$ & $4679.78\text{ ppm}$ & \textbf{$0.49\text{ ppm}$} & \textbf{$0.24\text{ ppm}$} \\
High Scalene $(10.0, 3.0, 1.0)$ & $1675.62\text{ ppm}$ & $3398.74\text{ ppm}$ & \textbf{$9.94\text{ ppm}$} & \textbf{$3.00\text{ ppm}$} \\
Prolate Spheroid $(10.0, 1.0, 1.0)$ & $8612.54\text{ ppm}$ & $7606.93\text{ ppm}$ & \textbf{$0.36\text{ ppm}$} & \textbf{$0.01\text{ ppm}$} \\
Oblate Spheroid $(10.0, 10.0, 1.0)$ & $7538.71\text{ ppm}$ & $10112.18\text{ ppm}$ & \textbf{$24.58\text{ ppm}$} & \textbf{$1.47\text{ ppm}$} \\
Extreme Prolate $(50.0, 1.0, 1.0)$ & $11385.12\text{ ppm}$ & $10235.32\text{ ppm}$ & \textbf{$0.42\text{ ppm}$} & \textbf{$0.01\text{ ppm}$} \\
Extreme Oblate $(50.0, 50.0, 1.0)$ & $7085.10\text{ ppm}$ & $10236.04\text{ ppm}$ & \textbf{$54.66\text{ ppm}$} & \textbf{$9.59\text{ ppm}$} \\
Flat Circular Disk $(100, 100, 0)$ & $6535.77\text{ ppm}$ & $9765.45\text{ ppm}$ & \textbf{$0.00\text{ ppm}$} & \textbf{$0.00\text{ ppm}$} \\
Flat Elliptical Disk $(100, 20, 0)$ & $6535.77\text{ ppm}$ & $9765.45\text{ ppm}$ & \textbf{$0.20\text{ ppm}$} & \textbf{$0.20\text{ ppm}$} \\
Extreme Triaxial Needle $(1000, 2.0, 1.0)$ & $6769.63\text{ ppm}$ & $5207.80\text{ ppm}$ & \textbf{$0.61\text{ ppm}$} & \textbf{$0.11\text{ ppm}$} \\
\bottomrule
\end{tabular}
\end{table}

\begin{thebibliography}{99}

\bibitem{Ivory1796}
J. Ivory, \emph{A new series for the rectification of the ellipsis; together with some observations on the evolution of the formula $(a^2 + b^2 - 2ab\cos\phi)^n$}, Transactions of the Royal Society of Edinburgh \textbf{4} (1796), 177--190.

\bibitem{Bessel1825}
F. W. Bessel, \emph{Ueber die Berechnung der geod\"atischen Linien}, Astronomische Nachrichten \textbf{4} (1825), no. 86, 241--254.

\bibitem{Ramanujan1914}
S. Ramanujan, \emph{Modular equations and approximations to $\pi$}, Quarterly Journal of Pure and Applied Mathematics \textbf{45} (1914), 350--372.

\bibitem{AbramowitzStegun1964}
M. Abramowitz and I. A. Stegun (eds.), \emph{Handbook of Mathematical Functions with Formulas, Graphs, and Mathematical Tables}, National Bureau of Standards, Applied Mathematics Series \textbf{55}, Washington, D.C., 1964.

\bibitem{JacobsenWaadeland1985}
L. Jacobsen and H. Waadeland, \emph{Some useful formulas involvingtails of continued fractions}, Lecture Notes in Mathematics \textbf{1199}, Springer, Berlin, 1986, 79--89.

\bibitem{AlmkvistBerndt1988}
G. Almkvist and B. C. Berndt, \emph{Gauss, Landen, Ramanujan, the arithmetic-geometric mean, ellipses, $\pi$, and the Ladies Diary}, The American Mathematical Monthly \textbf{95} (1988), no. 7, 585--608.

\bibitem{LorentzenWaadeland1992}
L. Lorentzen and H. Waadeland, \emph{Continued Fractions with Applications}, Studies in Computational Mathematics \textbf{3}, North-Holland, Amsterdam, 1992.

\bibitem{Carlson1995}
B. C. Carlson, \emph{Numerical computation of real or complex elliptic integrals}, Numerical Algorithms \textbf{10} (1995), no. 1, 13--26.

\bibitem{Cantrell2004}
D. W. Cantrell, \emph{Approximations for the perimeter of an ellipse and surface area of an ellipsoid}, \texttt{sci.math} archives and \emph{Numericana} compendium, May 2004.

\bibitem{Thomsen2004}
K. Thomsen, \emph{Symmetric power-mean approximation for the surface area of a scalene ellipsoid}, \emph{Numericana} compendium (ed. G. P. Michon), April 2004.

\bibitem{Villarino2005}
M. B. Villarino, \emph{Ramanujan's perimeter of an ellipse}, \texttt{arXiv:math/0506384} [math.CA], 2005 (published in J. Inequal. Pure Appl. Math. \textbf{9} (2008), Art. 70).

\bibitem{Sykora2005}
S. S\'ykora, \emph{Approximations of ellipse perimeters and of the complete elliptic integral $E(x)$: Review of known formulas}, Stan's Library, Vol. I, Castano Primo, Italy, 2005, \texttt{DOI:10.3247/SL1Math05.004}.

\bibitem{Sykora2007}
S. S\'ykora, \emph{Advances in approximations of ellipse perimeters and of the complete elliptic integral $E(x)$}, Stan's Library, Vol. II, Castano Primo, Italy, March 2007, \texttt{DOI:10.3247/SL2Math07.001}.

\bibitem{MichonNumericana}
G. P. Michon, \emph{Perimeter of an Ellipse \& Surface Area of an Ellipsoid}, Numericana Mathematical Tables (\texttt{http://www.numericana.com/answer/ellipse.htm}).

\end{thebibliography}

\end{document}
